% ID: SEL-D0040-OA-E01
% ESTADO: validada. S041, cotejada con examen y criterios.
\textbf{a)} Sacamos $X$ como factor por la derecha:
$(A+B)X=C$. Como $A+B=\begin{pmatrix}1&0\\2&1\end{pmatrix}$,
la primera fila de $X$ coincide con la primera de $C$.
Al restar dos veces esa fila de la segunda de $C$, resulta
\[
X=\boxed{\begin{pmatrix}2&1\\-1&0\end{pmatrix}}.
\]
El producto recupera las filas $(2,1)$ y $(3,2)$ de $C$.

\textbf{b)} Al multiplicar, $A^2=-I_2$, de modo que
$A^4=(-I_2)^2=\boxed{I_2}$. Como $80=4(20)$,
$A^{80}=(A^4)^{20}=\boxed{I_2}$.
